\chapter{大一时间线概览}

在大学生活这款开放式沙盒游戏里，我们总会在各种时间阶段遇到各种“事件”。本章试图将第一年的“事件”连缀成章，绘制大一生活的一张草图，帮助大家更快速地适应、更从容地规划丰富多彩的校园生活。也正因北大的生活过于多彩、选择过于多元，本章一定是平平无奇、挂一漏万的。因此，本章不是大一年级的“通关攻略”，而更应视作一张标注了NPC大致位置的参考地图。

\vspace{10pt}

此外，本章会加入大量索引，链接到前文内容。因此，也可以在茫茫信息涌到面前时，将本章作为一个大型、成文的目录，以便按需阅读全书内容。

\section{第一学期（秋季学期·开学前）}

\vspace{10pt}

\begin{center}
\begin{tikzpicture}[yscale=1.0, transform shape]
    \draw[chapc, line width=2pt] (0,0.35) -- (0,-13.55);

    % --- 1~4: 2倍间距 (4.4cm) ---

    % --- 1 ---
    \fill[chaph] (0,0) circle (4.5pt);
    \node[anchor=east, font=\small\bfseries, text=chapb, text width=3cm, align=right] at (-0.5,0) {8月下旬--9月上旬};
    \node[anchor=west, text width=11.3cm] at (0.5,0) {\textbf{英语分级考试}\\[2pt]{\footnotesize\color{chapc} 测试一般发生在军训后、开学前，题型涉及听力选择、阅读理解、词汇、语法、完形填空等，都是选择题，判卷速度很快(x)。正如3.3.4节所述，大家以最自然的状态参与分级测试就是最好的，适合的才是最好的。当然，如果你希望争取免修，或感觉直接去考反映不出自己真实水平，或就是莫名地不放心，可以找来各类英语考试的题目做一做。然而，笔者认为，为此刷题是不值当甚至不理性的。}};

    % --- 2 ---
    \fill[chaph] (0,-4.4) circle (4.5pt);
    \node[anchor=east, font=\small\bfseries, text=chapb, text width=3cm, align=right] at (-0.5,-4.4) {9月上旬};
    \node[anchor=west, text width=11.3cm] at (0.5,-4.4) {\textbf{新生训练营、新生定向越野}\\[2pt]{\footnotesize\color{chapc} 新生训练营帮助大家熟悉校园，做一些破冰活动，同时会有开学第一跑，可能会有神秘嘉宾哦！在定向越野中，你将和宿舍的同学一起在园内四处奔跑，完成一些有趣的游戏，继续熟悉校园。}};

    % --- 3 ---
    \fill[chaph] (0,-8.8) circle (4.5pt);
    \node[anchor=east, font=\small\bfseries, text=chapb, text width=3cm, align=right] at (-0.5,-8.8) {开学前};
    \node[anchor=west, text width=11.3cm] at (0.5,-8.8) {\textbf{选课}\\[2pt]{\footnotesize\color{chapc} 选课是门大学问，信息浩如烟海。详见2.3-2.5，也建议大家善用搜索。}};

    % --- 4 ---
    \fill[chaph] (0,-13.2) circle (4.5pt);
    \node[anchor=east, font=\small\bfseries, text=chapb, text width=3cm, align=right] at (-0.5,-13.2) {第1、2周};
    \node[anchor=west, text width=11.3cm] at (0.5,-13.2) {\textbf{第一次学业导师交流}\\[2pt]{\footnotesize\color{chapc} 笔者建议大家珍视同学业导师的每一次谈话。可以尝试在日常记录下自己遇到的、或大或小的困惑与问题，整合起来便于提问。当然，学业导师多由青年教师担任，因此大家不必焦虑，更不必以为学业导师全知全能，只要自然、愉快地开展沟通就好。详见3.6节}};


\end{tikzpicture}
\end{center}

\section{第一学期（秋季学期·期中前）}

\vspace{10pt}

\begin{center}
\begin{tikzpicture}[yscale=1.0, transform shape]
    \draw[chapc, line width=2pt] (0,0.35) -- (0,-16.85);

    % --- 5~8: 1.5倍间距 (3.3cm) ---

    % --- 5 ---
    \fill[chaph] (0,0) circle (4.5pt);
    \node[anchor=east, font=\small\bfseries, text=chapb, text width=3cm, align=right] at (-0.5,0) {第1、2周};
    \node[anchor=west, text width=11.3cm] at (0.5,0) {\textbf{工学院图书漂流活动}\\[2pt]{\footnotesize\color{chapc} 一般会在开学后的某个周末举行，让大家获取学长学姐的教材和参考书，节约购书成本。}};

    % --- 6 ---
    \fill[chaph] (0,-3.3) circle (4.5pt);
    \node[anchor=east, font=\small\bfseries, text=chapb, text width=3cm, align=right] at (-0.5,-3.3) {约第二周周末};
    \node[anchor=west, text width=11.3cm] at (0.5,-3.3) {\textbf{``百团大战''社团招新}\\[2pt]{\footnotesize\color{chapc} 上百社团同时招新，所以叫"百团大战"，在新太阳学生中心与百讲区域举办，是体验社团（整活）文化的绝佳机会，见4.3节。}};

    % --- 7 ---
    \fill[chaph] (0,-6.6) circle (4.5pt);
    \node[anchor=east, font=\small\bfseries, text=chapb, text width=3cm, align=right] at (-0.5,-6.6) {9月下旬};
    \node[anchor=west, text width=11.3cm] at (0.5,-6.6) {\textbf{秋季运动会}\\[2pt]{\footnotesize\color{chapc} 北大在春夏秋冬都会举办运动会，侧重的项目有所不同，其中秋季运动会也是新生趣味运动会。感兴趣的同学可以实时关注相关公众号。}};

    % --- 8 ---
    \fill[chaph] (0,-9.9) circle (4.5pt);
    \node[anchor=east, font=\small\bfseries, text=chapb, text width=3cm, align=right] at (-0.5,-9.9) {中秋节附近};
    \node[anchor=west, text width=11.3cm] at (0.5,-9.9) {\textbf{中秋晚会}\\[2pt]{\footnotesize\color{chapc} 每逢佳节，学校和各院系一般会分别举行晚会，工院一般会举行中秋晚会、新年晚会、毕业歌会等，随着工院的不断壮大，我们请到的嘉宾也越来越重磅。晚会不仅是一种视听盛宴，还能抽奖获得绝佳礼品哦！}};

    % --- 9 ---
    \fill[chaph] (0,-13.2) circle (4.5pt);
    \node[anchor=east, font=\small\bfseries, text=chapb, text width=3cm, align=right] at (-0.5,-13.2) {10月--12月上旬};
    \node[anchor=west, text width=11.3cm] at (0.5,-13.2) {\textbf{爱乐传习排练}\\[2pt]{\footnotesize\color{chapc} 爱乐传习是思政实践课程的一个选项（另一个是志愿服务），在排练前会先进行试音和分声部，同学们一般会在一周的固定几个时间段前往新奥大楼进行排练，``一二·九''歌会的排练准备阶段，持续约一个多月。}};

    % --- 10 ---
    \fill[chaph] (0,-16.5) circle (4.5pt);
    \node[anchor=east, font=\small\bfseries, text=chapb, text width=3cm, align=right] at (-0.5,-16.5) {第8周左右};
    \node[anchor=west, text width=11.3cm] at (0.5,-16.5) {\textbf{期中季}\\[2pt]{\footnotesize\color{chapc} 期中考试来临。同时可进行中期退课、自选PF（Pass/Fail）等操作，详见2.5节。注意紧密关注相关通知和截止日期！期中考试往往会使部分同学产生巨大的心理落差，希望大家调整好心态，调整好学习方法，不要太拘泥于考试结果之中}};
\end{tikzpicture}
\end{center}

\section{第一学期（秋季学期·期末）}

\vspace{10pt}

\begin{center}
\begin{tikzpicture}[yscale=1.0, transform shape]
    \draw[chapc, line width=2pt] (0,0.35) -- (0,-20.15);

    % --- 11 ---
    \fill[chaph] (0,0) circle (4.5pt);
    \node[anchor=east, font=\small\bfseries, text=chapb, text width=3cm, align=right] at (-0.5,0) {10/11月};
    \node[anchor=west, text width=11.3cm] at (0.5,0) {\textbf{新生舞会}\\[2pt]{\footnotesize\color{chapc} 多院系联谊活动，认识其他院系的同学，体验大学生活的多彩一面。现场会有舞培、表演和丰富的茶歇，可以自带舞伴，也可以填写问卷进行匹配，注意关注报名信息。}};


    % --- 12 ---
    \fill[chaph] (0,-3.3) circle (4.5pt);
    \node[anchor=east, font=\small\bfseries, text=chapb, text width=3cm, align=right] at (-0.5,-3.3) {11月--12月};
    \node[anchor=west, text width=11.3cm] at (0.5,-3.3) {\textbf{工学嘉年华}\\[2pt]{\footnotesize\color{chapc} 《现代工学通论》配套活动。是了解各专业方向、与教授和学长学姐交流的绝佳机会，每专业一场，分为讲座介绍和实验室实地参观，让同学们建立对各个专业的大概认知。打卡若干场活动才能完成课程要求！}};

    % --- 13 ---
    \fill[chaph] (0,-6.6) circle (4.5pt);
    \node[anchor=east, font=\small\bfseries, text=chapb, text width=3cm, align=right] at (-0.5,-6.6) {12月上旬};
    \node[anchor=west, text width=11.3cm] at (0.5,-6.6) {\textbf{``一二·九''歌会}\\[2pt]{\footnotesize\color{chapc} 全校性新生合唱比赛，纪念``一二·九''运动。爱乐传习排练的最终呈现，一般在百周年纪念讲堂举行。}};

    % --- 14 ---
    \fill[chaph] (0,-9.9) circle (4.5pt);
    \node[anchor=east, font=\small\bfseries, text=chapb, text width=3cm, align=right] at (-0.5,-9.9) {期末周前};
    \node[anchor=west, text width=11.3cm] at (0.5,-9.9) {\textbf{四六级考试}\\[2pt]{\footnotesize\color{chapc} 全国大学英语四、六级考试，大一上可以参加四级考试，由于大一时英语内容的遗忘程度较低，建议尽早参加。同时，四六级考试一般安排在期末周前，需要平衡好各场考试的复习时间，不建议完全裸考。}};

    % --- 15 ---
    \fill[chaph] (0,-13.2) circle (4.5pt);
    \node[anchor=east, font=\small\bfseries, text=chapb, text width=3cm, align=right] at (-0.5,-13.2) {12月底};
    \node[anchor=west, text width=11.3cm] at (0.5,-13.2) {\textbf{新年晚会}\\[2pt]{\footnotesize\color{chapc} 与中秋晚会类似，各院系和学校层面均有举办。}};

    % --- 16 ---
    \fill[chaph] (0,-16.5) circle (4.5pt);
    \node[anchor=east, font=\small\bfseries, text=chapb, text width=3cm, align=right] at (-0.5,-16.5) {12月底--次年1月初};
    \node[anchor=west, text width=11.3cm] at (0.5,-16.5) {\textbf{期末周}\\[2pt]{\footnotesize\color{chapc} 停课复习考试，提前规划复习，合理分配时间（尤其是考试十分密集的情况），注意通识课一般在学期内最后一个行课周进行考试，而不是期末周。注意图书馆开放时间、浴室开放时间在期末周均有所延长}};

    % --- 17 ---
    \fill[chaph] (0,-19.8) circle (4.5pt);
    \node[anchor=east, font=\small\bfseries, text=chapb, text width=3cm, align=right] at (-0.5,-19.8) {1月};
    \node[anchor=west, text width=11.3cm] at (0.5,-19.8) {\textbf{冬季运动会}\\[2pt]{\footnotesize\color{chapc} }};
\end{tikzpicture}
\end{center}

\section{第二学期（春季学期）}

\vspace{10pt}

\begin{center}
\begin{tikzpicture}[yscale=1.0, transform shape]
    \draw[chapc, line width=2pt] (0,0.35) -- (0,-22.35);

    % --- 18~27: 2.2cm间距（不变） ---

    % --- 18 ---
    \fill[chaph] (0,0) circle (4.5pt);
    \node[anchor=east, font=\small\bfseries, text=chapb, text width=3cm, align=right] at (-0.5,0) {约第二周周末};
    \node[anchor=west, text width=11.3cm] at (0.5,0) {\textbf{百团大战（春季）}\\[2pt]{\footnotesize\color{chapc} 春季学期社团招新}};

    % --- 19 ---
    \fill[chaph] (0,-2.2) circle (4.5pt);
    \node[anchor=east, font=\small\bfseries, text=chapb, text width=3cm, align=right] at (-0.5,-2.2) {3月};
    \node[anchor=west, text width=11.3cm] at (0.5,-2.2) {\textbf{工学文化节}\\[2pt]{\footnotesize\color{chapc} 在百周年纪念讲堂前举办，有许多有趣的科普内容或小游戏，如机器狗表演、选课主题桌游等，还可以兑换礼品，欢迎前来捧场！}};

    % --- 20 ---
    \fill[chaph] (0,-4.4) circle (4.5pt);
    \node[anchor=east, font=\small\bfseries, text=chapb, text width=3cm, align=right] at (-0.5,-4.4) {3月底};
    \node[anchor=west, text width=11.3cm] at (0.5,-4.4) {\textbf{十佳歌手决赛}\\[2pt]{\footnotesize\color{chapc} 校园歌手大赛决赛，在邱德拔体育馆举行，是北大最具人气的文艺活动之一，如果想要位置不错的票，记得提前去排队！}};

    % --- 21 ---
    \fill[chaph] (0,-6.6) circle (4.5pt);
    \node[anchor=east, font=\small\bfseries, text=chapb, text width=3cm, align=right] at (-0.5,-6.6) {4月};
    \node[anchor=west, text width=11.3cm] at (0.5,-6.6) {\textbf{工学院积分大赛}\\[2pt]{\footnotesize\color{chapc} 工学院特色活动，分为初赛、复赛和决赛，奖品丰富，题目具有难度层次，具体情况可见往届推送。欢迎来检验自己的积分能力！}};

    % --- 22 ---
    \fill[chaph] (0,-8.8) circle (4.5pt);
    \node[anchor=east, font=\small\bfseries, text=chapb, text width=3cm, align=right] at (-0.5,-8.8) {4月};
    \node[anchor=west, text width=11.3cm] at (0.5,-8.8) {\textbf{春季运动会}\\[2pt]{\footnotesize\color{chapc} }};

    % --- 23 ---
    \fill[chaph] (0,-11.0) circle (4.5pt);
    \node[anchor=east, font=\small\bfseries, text=chapb, text width=3cm, align=right] at (-0.5,-11.0) {第8周左右};
    \node[anchor=west, text width=11.3cm] at (0.5,-11.0) {\textbf{期中季}\\[2pt]{\footnotesize\color{chapc} }};

    % --- 24 ---
    \fill[chaph] (0,-13.2) circle (4.5pt);
    \node[anchor=east, font=\small\bfseries, text=chapb, text width=3cm, align=right] at (-0.5,-13.2) {4月--5月};
    \node[anchor=west, text width=11.3cm] at (0.5,-13.2) {\textbf{双专业报名、转专业}\\[2pt]{\footnotesize\color{chapc} 重要的学业决策窗口，学院和学校都会提前举办相关分享讲座，共享信息，有意向的同学可以关注或咨询教务老师。详见3.4和3.5节，笔者在这里强调提前搜索相关信息的重要性，比如在树洞搜索关键词阅读往年的申请情况}};

    % --- 25 ---
    \fill[chaph] (0,-15.4) circle (4.5pt);
    \node[anchor=east, font=\small\bfseries, text=chapb, text width=3cm, align=right] at (-0.5,-15.4) {期末周前};
    \node[anchor=west, text width=11.3cm] at (0.5,-15.4) {\textbf{四六级考试}\\[2pt]{\footnotesize\color{chapc} 对大多数同学，一般是参加六级考试}};

    % --- 26 ---
    \fill[chaph] (0,-17.6) circle (4.5pt);
    \node[anchor=east, font=\small\bfseries, text=chapb, text width=3cm, align=right] at (-0.5,-17.6) {6月中旬};
    \node[anchor=west, text width=11.3cm] at (0.5,-17.6) {\textbf{期末周}\\[2pt]{\footnotesize\color{chapc} 期末考试之后便是专业方向确认与分流，注意学院通知。}};

    % --- 27 ---
    \fill[chaph] (0,-19.8) circle (4.5pt);
    \node[anchor=east, font=\small\bfseries, text=chapb, text width=3cm, align=right] at (-0.5,-19.8) {暑假};
    \node[anchor=west, text width=11.3cm] at (0.5,-19.8) {\textbf{思政实践}\\[2pt]{\footnotesize\color{chapc} 暑期社会实践课程（必修），一般是探访实践地的相关政府、企业、文化馆等，了解当地相关产业的发展。}};

    % --- 28 ---
    \fill[chaph] (0,-22.0) circle (4.5pt);
    \node[anchor=east, font=\small\bfseries, text=chapb, text width=3cm, align=right] at (-0.5,-22.0) {暑假};
    \node[anchor=west, text width=11.3cm] at (0.5,-22.0) {\textbf{海外暑校}\\[2pt]{\footnotesize\color{chapc} 2-3月份可以开始关注暑期研学课程（海外暑校）。}};
 
\end{tikzpicture}
\end{center}

\vspace{10pt}

以上时间线涵盖了北大工学院大一学年的主要事件节点。希望对同学们把握全年节奏有所帮助。当然，大学中还有很多未被列出的精彩——课堂上的灵光一现、图书馆里的专注时光、未名湖畔的漫步、与好友的深夜畅谈……这些无法被时间线标注的时刻，往往才是大学生活中最珍贵的部分。愿你在参考时间线的同时，也留出空间去发现属于自己的节奏与风景。
